\documentclass[12pt,a4paper]{report}

\usepackage[T1]{fontenc}
\usepackage[utf8]{inputenc}
\usepackage[french]{babel}
\usepackage{newtxtext,newtxmath}
\usepackage{microtype}
\usepackage[a4paper,top=2.5cm,bottom=2.5cm,left=2.7cm,right=2.4cm]{geometry}
\usepackage{setspace}
\onehalfspacing
\setlength{\parindent}{0pt}
\setlength{\parskip}{6pt}
\usepackage{graphicx}
\usepackage{float}
\usepackage{array}
\usepackage{booktabs}
\usepackage{tabularx}
\usepackage{longtable}
\usepackage{multirow}
\usepackage[table]{xcolor}
\usepackage{colortbl}
\usepackage{caption}
\usepackage{enumitem}
\usepackage{tikz}
\usetikzlibrary{positioning,arrows.meta,shapes.geometric,fit}
\usepackage{mdframed}
\usepackage[most]{tcolorbox}
\usepackage{placeins}
\usepackage{hyperref}
\usepackage{fancyhdr}
\usepackage{titlesec}

\definecolor{mainblue}{RGB}{0,102,204}
\definecolor{darkblue}{RGB}{0,55,110}
\definecolor{lightblue}{RGB}{230,242,255}
\definecolor{softblue}{RGB}{242,248,255}
\definecolor{graytitle}{RGB}{70,70,70}
\definecolor{rowgray}{RGB}{245,247,250}
\definecolor{successgreen}{RGB}{39,174,96}
\definecolor{warningorange}{RGB}{230,126,34}
\definecolor{dangerred}{RGB}{192,57,43}
\definecolor{purple}{RGB}{142,68,173}

\hypersetup{
  colorlinks=true,
  linkcolor=mainblue,
  urlcolor=mainblue,
  citecolor=mainblue,
  pdftitle={Sprint 3 - Supervision, Reporting, Messagerie et Audit},
  pdfauthor={HandiTalents}
}

\setcounter{secnumdepth}{3}
\setcounter{tocdepth}{3}
\pagestyle{fancy}
\fancyhf{}
\fancyhead[L]{\textit{HandiTalents}}
\fancyhead[R]{\textit{Sprint 3 : Supervision, Reporting, Messagerie et Audit}}
\fancyfoot[C]{\thepage}
\renewcommand{\headrulewidth}{0.4pt}
\renewcommand{\footrulewidth}{0pt}
\setlength{\headheight}{14.5pt}

\titleformat{\chapter}[display]
{\bfseries\Huge\color{mainblue}}
{\filleft\Large\chaptername~\thechapter}
{1ex}
{\titlerule\vspace{1.5ex}\filright}
[\vspace{1ex}\titlerule]
\titleformat{\section}{\Large\bfseries\color{mainblue}}{\thesection}{0.8em}{}
\titleformat{\subsection}{\large\bfseries\color{darkblue}}{\thesubsection}{0.8em}{}
\titleformat{\subsubsection}{\normalsize\bfseries\color{graytitle}}{\thesubsubsection}{0.8em}{}

\captionsetup{font=small,labelfont=bf,textfont=it,justification=centering,skip=8pt}
\setlist[itemize]{leftmargin=1.5em,itemsep=3pt,topsep=3pt}
\setlist[enumerate]{leftmargin=1.8em,itemsep=3pt,topsep=3pt}

\newtcolorbox{analysisbox}[1][]{enhanced,breakable,colback=white,colframe=mainblue!45,
  coltitle=white,fonttitle=\bfseries,title=#1,boxrule=0.6pt,arc=2mm,
  left=7pt,right=7pt,top=6pt,bottom=6pt}

\newcommand{\safeimage}[2]{%
  \IfFileExists{#2}{%
    \includegraphics[width=#1]{#2}%
  }{%
    \fbox{%
      \parbox[c][5cm][c]{#1}{\centering\small Image manquante\\\texttt{\detokenize{#2}}}%
    }%
  }%
}

\newcommand{\framedimage}[2]{%
\begin{mdframed}[linecolor=mainblue!40,linewidth=0.6pt,roundcorner=4pt,
  backgroundcolor=white,innertopmargin=5pt,innerbottommargin=5pt,
  innerleftmargin=5pt,innerrightmargin=5pt]
  \centering\safeimage{#1}{#2}
\end{mdframed}}

\newcommand{\bigscreen}[4]{%
\begin{figure}[H]
  \centering
  \framedimage{#1}{#2}
  \caption{#3}
  \label{#4}
\end{figure}
\FloatBarrier
}

\newcommand{\badgePending}{\colorbox{warningorange!18}{\textcolor{warningorange}{\textbf{Submitted}}}}
\newcommand{\badgeValidated}{\colorbox{successgreen!15}{\textcolor{successgreen}{\textbf{Validated}}}}
\newcommand{\badgeRejected}{\colorbox{dangerred!15}{\textcolor{dangerred}{\textbf{Rejected}}}}
\newcommand{\badgeActive}{\colorbox{mainblue!15}{\textcolor{mainblue}{\textbf{Active}}}}
\newcommand{\badgeSuspended}{\colorbox{dangerred!15}{\textcolor{dangerred}{\textbf{Suspended}}}}
\newcommand{\badgeArchived}{\colorbox{graytitle!15}{\textcolor{graytitle}{\textbf{Archived}}}}

\newcolumntype{Y}{>{\raggedright\arraybackslash}X}
\newcolumntype{L}[1]{>{\raggedright\arraybackslash}p{#1}}

\begin{document}

\begin{titlepage}
  \centering
  \vspace*{3.4cm}
  {\Huge\bfseries\color{mainblue} Sprint 3 \\[0.4cm]
  Supervision inclusive, Reporting, Messagerie temps réel et Audit\par}
  \vspace{1.2cm}
  \rule{12cm}{0.7pt}
  \vspace{1.2cm}
  {\Large Mise en place des modules de gouvernance, de conformité,
  de communication et de traçabilité de HandiTalents.\par}
  \vfill
\end{titlepage}

\chapter{Sprint 3 : Supervision inclusive, Reporting, Messagerie et Audit}

\section{Introduction}

Après l'authentification et la gestion des profils au Sprint~1, puis les tests Soft Skills,
les candidatures et les entretiens au Sprint~2, le Sprint~3 consolide HandiTalents par des
fonctionnalités transverses. Ce sprint apporte une vraie couche de gouvernance : supervision
inclusive pour les acteurs institutionnels, reporting réglementaire pour les entreprises,
messagerie temps réel pour la collaboration, gestion avancée des utilisateurs pour
l'administrateur et audit complet pour la traçabilité.

L'objectif n'est pas seulement d'ajouter des écrans. Il s'agit de rendre la plateforme plus
pilotable, plus vérifiable et plus crédible dans un contexte de recrutement inclusif.
Chaque fonctionnalité a été pensée avec un contrôle d'accès strict, des données filtrées
selon le rôle et une exposition claire des workflows métier.

\section{Phase de préparation}

\subsection{Objectif}

L'objectif du Sprint~3 est de transformer HandiTalents en une plateforme de supervision
et de conformité à part entière.

\textbf{Du point de vue technique}, ce sprint vise à :
\begin{itemize}
  \item exposer des endpoints sécurisés pour la supervision, le reporting entreprise,
  la messagerie, la gestion des utilisateurs, les membres d'équipe et l'audit ;
  \item structurer les parcours autour du RBAC et du filtrage territorial ;
  \item stocker et relier proprement les rapports, conversations, messages, membres et
  événements d'audit ;
  \item diffuser les messages en temps réel via SSE ;
  \item enregistrer les opérations sensibles dans \texttt{audit\_actions\_admin}.
\end{itemize}

\textbf{Du point de vue fonctionnel}, ce sprint vise à :
\begin{itemize}
  \item permettre à l'inspecteur territorial et à l'ANETI de suivre la plateforme avec
  des tableaux de bord adaptés à leur périmètre ;
  \item permettre aux entreprises de créer, suivre et transmettre leurs rapports de
  conformité ;
  \item permettre aux utilisateurs authentifiés d'échanger sans rechargement de page ;
  \item permettre à l'administrateur de gérer les comptes avec audit et réinitialisation ;
  \item permettre à l'entreprise d'organiser ses membres internes.
\end{itemize}

\subsection{Backlog du Sprint~3}

Le tableau~\ref{tab:sprint3_backlog} présente les user stories retenues pour ce sprint.
La définition de terminé précise les comportements attendus et les critères de validation.

\begin{longtable}{|>{\centering\arraybackslash}p{2.7cm}|>{\centering\arraybackslash}p{1.3cm}|>{\raggedright\arraybackslash}p{4.8cm}|>{\raggedright\arraybackslash}p{4.8cm}|>{\centering\arraybackslash}p{1.2cm}|}
\caption{Backlog du Sprint~3 : Supervision, Reporting, Messagerie et Audit}
\label{tab:sprint3_backlog}\\
\hline
\rowcolor{mainblue!15}
\textbf{Thème} & \textbf{US} & \textbf{User Story} & \textbf{Definition of Done} & \textbf{Est.} \\
\hline
\endfirsthead
\hline
\rowcolor{mainblue!15}
\textbf{Thème} & \textbf{US} & \textbf{User Story} & \textbf{Definition of Done} & \textbf{Est.} \\
\hline
\endhead

\multirow{4}{2.7cm}{\centering\textbf{Supervision\\inclusive}}
& US-01 & En tant qu'inspecteur ou agent ANETI, je veux consulter un dashboard avec KPIs afin de suivre le recrutement inclusif. & Les indicateurs sont visibles selon le rôle : vision territoriale pour l'inspecteur et nationale pour l'ANETI. & 3j \\
\cline{2-5}
& US-02 & En tant qu'acteur de supervision, je veux visualiser le pipeline de recrutement afin d'identifier les points de friction. & Les volumes par étape sont lisibles et filtrables par période, entreprise et territoire. & 2j \\
\cline{2-5}
& US-03 & En tant qu'acteur de supervision, je veux consulter les rapports de conformité afin de les valider ou les rejeter. & La liste et le détail des rapports sont accessibles, le rejet impose une recommandation. & 3j \\
\cline{2-5}
& US-04 & En tant qu'acteur de supervision, je veux exporter les données filtrées afin de produire des analyses externes. & L'export CSV/Excel applique les filtres actifs et les règles RBAC. & 1.5j \\
\hline

\multirow{2}{2.7cm}{\centering\textbf{Reporting\\entreprise}}
& US-05 & En tant qu'entreprise, je veux consulter mon hub de conformité afin de suivre ma situation réglementaire. & Les rapports, indicateurs et canaux de publication sont visibles depuis l'espace entreprise. & 2j \\
\cline{2-5}
& US-06 & En tant qu'entreprise, je veux créer et soumettre un rapport de conformité. & Le rapport est enregistré avec le statut \texttt{submitted} et devient exploitable côté supervision. & 3j \\
\hline

\multirow{2}{2.7cm}{\centering\textbf{Gestion\\utilisateurs}}
& US-07 & En tant qu'administrateur, je veux gérer les utilisateurs de la plateforme. & La liste, la recherche, la création, la modification, la suppression, les statuts et l'export sont disponibles. & 4j \\
\cline{2-5}
& US-08 & En tant qu'administrateur, je veux réinitialiser le mot de passe d'un utilisateur. & Un accès temporaire est généré et l'opération est enregistrée dans l'audit. & 1j \\
\hline

\multirow{2}{2.7cm}{\centering\textbf{Messagerie}}
& US-09 & En tant qu'utilisateur authentifié, je veux créer une conversation et envoyer des messages. & Les conversations sont persistées, les participants contrôlés et l'historique est paginé. & 2.5j \\
\cline{2-5}
& US-10 & En tant qu'utilisateur, je veux recevoir les messages en temps réel. & Le flux SSE diffuse immédiatement les nouveaux messages aux participants connectés. & 2.5j \\
\hline

\multirow{1}{2.7cm}{\centering\textbf{Equipe\\entreprise}}
& US-11 & En tant qu'entreprise, je veux gérer les membres de mon équipe. & L'ajout, la modification et la suppression des membres fonctionnent dans le périmètre de l'entreprise. & 2j \\
\hline

\multirow{2}{2.7cm}{\centering\textbf{Audit}}
& US-12 & En tant que système, je veux tracer les actions sensibles. & Les valeurs avant/après, l'acteur, la cible, le type d'action et la date sont enregistrés. & 2j \\
\cline{2-5}
& US-13 & En tant qu'administrateur, je veux consulter l'historique des actions. & L'historique d'audit est consultable avec filtres par type et période. & 1j \\
\hline
\rowcolor{mainblue!10}
\multicolumn{4}{|r|}{\textbf{Total estimé}} & \textbf{29.5j} \\
\hline
\end{longtable}

\section{Phase de mise en \oe uvre}

\subsection{Modélisation fonctionnelle}

\subsubsection{Diagramme de cas d'utilisation du Sprint~3}

Le diagramme de cas d'utilisation présente les cinq familles d'acteurs du sprint :
\textbf{inspecteur territorial}, \textbf{agent ANETI}, \textbf{entreprise},
\textbf{administrateur} et \textbf{utilisateur authentifié}. Les cas sensibles passent
par l'authentification et les droits sont ajustés selon le rôle. L'inspecteur opère sur
son territoire, l'ANETI sur l'ensemble de la plateforme et l'administrateur sur les
comptes, l'audit et la gouvernance.

\bigscreen{0.95\textwidth}{images/sprint3_usecase.png}%
{Diagramme de cas d'utilisation global — Sprint~3}{fig:s3_usecase}

\small
\setlength{\tabcolsep}{4pt}
\renewcommand{\arraystretch}{1.15}
\begin{longtable}{|L{1.3cm}|L{3.9cm}|L{3.1cm}|L{6.1cm}|}
\caption{Description textuelle des cas d'utilisation — Sprint~3}
\label{tab:s3_uc_description}\\
\hline
\rowcolor{mainblue!15}
\textbf{ID} & \textbf{Cas d'utilisation} & \textbf{Acteur(s)} & \textbf{Description} \\
\hline
\endfirsthead
\hline
\rowcolor{mainblue!15}
\textbf{ID} & \textbf{Cas d'utilisation} & \textbf{Acteur(s)} & \textbf{Description} \\
\hline
\endhead
UC-01 & Consulter le dashboard de supervision & Inspecteur, ANETI & Affichage des KPIs du recrutement inclusif avec un périmètre dépendant du rôle. \\
\hline
UC-02 & Visualiser le pipeline de recrutement & Inspecteur, ANETI & Suivi des candidatures, shortlists, entretiens et recrutements par entreprise. \\
\hline
UC-03 & Analyser les offres d'emploi & Inspecteur, ANETI & Consultation des indicateurs de performance et de conversion des offres. \\
\hline
UC-04 & Exporter les données & Inspecteur, ANETI & Export des données filtrées selon les droits et le territoire. \\
\hline
UC-05 & Gérer les rapports de conformité & Admin, Inspecteur, ANETI & Consultation, validation, rejet et ajout de recommandations. \\
\hline
UC-06 & Créer un rapport de conformité & Entreprise & Saisie et soumission d'un rapport réglementaire avec statut initial \texttt{submitted}. \\
\hline
UC-07 & Gérer les membres d'équipe & Entreprise & Administration des membres associés à l'espace entreprise. \\
\hline
UC-08 & Gérer les utilisateurs & Administrateur & CRUD, statuts, export, réinitialisation de mot de passe et recherche. \\
\hline
UC-09 & Consulter l'historique d'audit & Administrateur & Consultation des actions sensibles liées aux comptes et aux statuts. \\
\hline
UC-10 & Echanger par messagerie temps réel & Utilisateur authentifié & Création de conversations, envoi de messages et réception via SSE. \\
\hline
\end{longtable}
\normalsize

\subsubsection{Diagramme de classes du Sprint~3}

Le diagramme de classes enrichit le noyau métier avec les entités de conformité,
de messagerie, d'équipe et d'audit. Il permet d'identifier clairement les relations
entre les rapports, les recommandations, les conversations, les participants et les
actions administratives tracées.

\bigscreen{0.95\textwidth}{images/sprint3_class_diagram.png}%
{Diagramme de classes global — Sprint~3}{fig:s3_classes}

\begin{table}[H]
\centering
\caption{Entités principales du Sprint~3}
\label{tab:s3_entites}
\renewcommand{\arraystretch}{1.35}
\rowcolors{2}{rowgray}{white}
\begin{tabularx}{\textwidth}{|L{4cm}|Y|}
\hline
\rowcolor{mainblue!15}
\textbf{Entité} & \textbf{Rôle métier} \\
\hline
ComplianceReport & Rapport soumis par une entreprise et traité par les acteurs de supervision. \\
\hline
ComplianceRecommendation & Recommandation associée à un rapport, notamment en cas de rejet. \\
\hline
Conversation & Canal de discussion entre utilisateurs authentifiés. \\
\hline
ConversationParticipant & Association entre un utilisateur et une conversation pour sécuriser l'accès. \\
\hline
Message & Message envoyé dans une conversation et diffusé aux participants connectés. \\
\hline
AccountMember & Membre de l'équipe entreprise avec coordonnées et rôle interne. \\
\hline
AuditActionAdmin & Trace d'une action administrative sensible, avec valeurs avant et après. \\
\hline
\end{tabularx}
\end{table}

\subsection{Conception}

\subsubsection{Diagramme de séquence — Soumission et validation d'un rapport de conformité}

Le premier flux couvre la création d'un rapport par l'entreprise puis son traitement
par un acteur de supervision. La séquence montre une logique simple mais solide :
saisie, persistante, consultation, décision et notification.

\bigscreen{0.95\textwidth}{images/sprint3_sequence_report.png}%
{Diagramme de séquence — Soumission et validation d'un rapport de conformité}{fig:s3_seq_report}

\textbf{Déroulement du flux :}
\begin{enumerate}
  \item L'entreprise prépare son rapport et l'envoie via \texttt{POST /api/entreprise/reports-requests/reports}.
  \item Le backend persiste le rapport avec le statut \texttt{submitted} et conserve la structure de version.
  \item L'inspecteur ou l'ANETI récupère la liste via \texttt{GET /api/supervision/reports}.
  \item Le rapport est ensuite validé ou rejeté via \texttt{POST /api/supervision/reports/\{id\}/validate} ou \texttt{POST /api/supervision/reports/\{id\}/reject}.
  \item En cas de rejet, une recommandation devient obligatoire et l'entreprise est notifiée.
\end{enumerate}

\subsubsection{Diagramme de séquence — Messagerie temps réel}

Le second flux décrit la messagerie temps réel. L'objectif est d'éviter toute
expérience lourde côté utilisateur : la conversation est créée une seule fois, puis
les messages circulent via SSE sans rechargement.

\bigscreen{0.95\textwidth}{images/sprint3_sequence_chat.png}%
{Diagramme de séquence — Messagerie temps réel via SSE}{fig:s3_seq_chat}

\textbf{Déroulement du flux :}
\begin{enumerate}
  \item Un utilisateur crée une conversation et ajoute les participants autorisés.
  \item Le frontend ouvre un flux SSE sur \texttt{GET /api/chat/conversations/\{id\}/stream?token=JWT}.
  \item Un message est envoyé via \texttt{POST /api/chat/conversations/\{id\}/messages}.
  \item Le backend persiste le message puis le diffuse immédiatement aux abonnés connectés.
  \item Les autres participants voient le contenu apparaître sans rechargement.
\end{enumerate}

\subsubsection{Diagramme de séquence — Gestion des utilisateurs et audit}

Ce dernier flux est essentiel pour la gouvernance : toute opération sensible effectuée
par l'administrateur doit laisser une trace. Cela concerne les changements de statut,
les modifications de profil et la réinitialisation de mot de passe.

\bigscreen{0.95\textwidth}{images/sprint3_sequence_audit.png}%
{Diagramme de séquence — Gestion des utilisateurs et traçabilité audit}{fig:s3_seq_audit}

\textbf{Déroulement du flux :}
\begin{enumerate}
  \item L'administrateur modifie un utilisateur ou lance une réinitialisation de mot de passe.
  \item Le backend récupère d'abord l'état précédent pour pouvoir comparer avant/après.
  \item La modification est appliquée puis une entrée est enregistrée dans \texttt{audit\_actions\_admin}.
  \item Si nécessaire, un accès temporaire ou une notification est généré.
  \item L'action reste traçable dans l'historique exploitable par l'administrateur.
\end{enumerate}

\subsection{Architecture et organisation technique}

Le Sprint~3 s'appuie sur une organisation par domaines fonctionnels. Même si le projet
reste encore monolithique sur certains aspects, la séparation en routes et services permet
de préparer une architecture plus distribuée et plus robuste.

\begin{table}[H]
\centering
\caption{Modules techniques du Sprint~3}
\label{tab:s3_services}
\renewcommand{\arraystretch}{1.35}
\rowcolors{2}{rowgray}{white}
\begin{tabularx}{\textwidth}{|L{4.1cm}|L{4.8cm}|Y|}
\hline
\rowcolor{mainblue!15}
\textbf{Module} & \textbf{Endpoints principaux} & \textbf{Responsabilité} \\
\hline
Supervision & \texttt{/api/supervision/*} & KPIs, pipeline, rapports, offres, candidats et export. \\
\hline
Reporting entreprise & \texttt{/api/entreprise/reports-requests/*} & Contexte de conformité, historique et création de rapports. \\
\hline
Messagerie & \texttt{/api/chat/*} & Conversations, participants, messages et flux SSE. \\
\hline
Gestion utilisateurs & \texttt{/api/admin/utilisateurs/*} & CRUD, statuts, export, reset password et historique utilisateur. \\
\hline
Equipe entreprise & \texttt{/api/entreprises/membres/*} & Gestion des membres associés à une entreprise. \\
\hline
Audit & \texttt{audit\_actions\_admin} & Traçabilité des actions sensibles administratives. \\
\hline
\end{tabularx}
\end{table}

\subsection{Réalisation}

\subsubsection{Module 1 — Supervision inclusive}

Ce module est le plus visible pour les acteurs institutionnels. Il regroupe les
indicateurs de recrutement inclusif, la vue pipeline, la performance des offres et la
liste des rapports à traiter. L'objectif de l'interface n'est pas d'être décorative,
mais lisible et utile pour prendre des décisions.

\paragraph{Workflow du module}

\begin{table}[H]
\centering
\caption{Workflow du module Supervision inclusive}
\label{tab:s3_workflow_supervision}
\renewcommand{\arraystretch}{1.35}
\rowcolors{2}{rowgray}{white}
\begin{tabularx}{\textwidth}{|L{4cm}|Y|}
\hline
\rowcolor{mainblue!15}
\textbf{Étape} & \textbf{Description} \\
\hline
Chargement du tableau de bord & L'acteur de supervision accède aux indicateurs selon son rôle et son périmètre. \\
\hline
Lecture des KPIs & Le système affiche les volumes, taux et tendances utiles au pilotage. \\
\hline
Exploration du pipeline & Les candidatures, shortlists et entretiens sont visualisés par entreprise et par période. \\
\hline
Gestion des rapports & Les rapports de conformité sont ouverts, analysés, validés ou rejetés. \\
\hline
Export des données & Les données filtrées sont exportées pour les besoins de contrôle ou d'analyse. \\
\hline
\end{tabularx}
\end{table}

\paragraph{Interfaces réalisées}

Les écrans retenus pour ce module mettent en avant les fonctions les plus représentatives :
le \textbf{dashboard de supervision}, la \textbf{vue pipeline}, la \textbf{liste des rapports}
et l'\textbf{export des données}. Ces interfaces illustrent clairement le rôle du module
dans le pilotage de la plateforme.

\bigscreen{0.92\textwidth}{images/supervision_dashboard.png}%
{Dashboard de supervision inclusive}{fig:s3_ui_dashboard}

\bigscreen{0.92\textwidth}{images/supervision_reports.png}%
{Liste et traitement des rapports de conformité}{fig:s3_ui_reports}

\bigscreen{0.92\textwidth}{images/supervision_export.png}%
{Export de données de supervision}{fig:s3_ui_export}

\subsubsection{Module 2 — Reporting de conformité entreprise}

Ce module donne à l'entreprise une vue claire de sa situation réglementaire. L'espace
reporting doit rester simple : un hub de conformité, un formulaire de création, un
historique des versions et des canaux de publication explicites.

\paragraph{Workflow du module}

\begin{table}[H]
\centering
\caption{Workflow du module Reporting entreprise}
\label{tab:s3_workflow_reporting}
\renewcommand{\arraystretch}{1.35}
\rowcolors{2}{rowgray}{white}
\begin{tabularx}{\textwidth}{|L{4cm}|Y|}
\hline
\rowcolor{mainblue!15}
\textbf{Étape} & \textbf{Description} \\
\hline
Ouverture du hub & L'entreprise consulte ses indicateurs et ses rapports existants. \\
\hline
Préparation du rapport & Les données de conformité sont saisies dans un formulaire guidé. \\
\hline
Soumission & Le rapport est enregistré avec un statut interne de type \texttt{submitted}. \\
\hline
Suivi des statuts & Le rapport passe ensuite par les états \texttt{validated} ou \texttt{rejected}. \\
\hline
Publication & Les canaux de diffusion sont activés selon la décision de supervision. \\
\hline
\end{tabularx}
\end{table}

\paragraph{Interfaces réalisées}

Les interfaces retenues pour ce module sont le \textbf{hub de conformité}, le
\textbf{formulaire de création} et l'\textbf{historique des rapports}. Cet ensemble
présente la logique métier du reporting de manière lisible et progressive.

\bigscreen{0.92\textwidth}{images/company_reporting.png}%
{Hub de reporting de conformité}{fig:s3_ui_reporting}

\bigscreen{0.92\textwidth}{images/company_reporting_form.png}%
{Création d'un rapport de conformité}{fig:s3_ui_reporting_form}

\bigscreen{0.92\textwidth}{images/company_reporting_history.png}%
{Historique des rapports et canaux de publication}{fig:s3_ui_reporting_history}

\subsubsection{Module 3 — Gestion avancée des utilisateurs}

Le module de gestion des utilisateurs est central côté administrateur. Il doit rester
dense mais lisible : une liste paginée, des filtres clairs, une fiche de détail, un
formulaire de modification et les actions sensibles visibles sans surcharge visuelle.

\paragraph{Workflow du module}

\begin{table}[H]
\centering
\caption{Workflow du module Gestion avancée des utilisateurs}
\label{tab:s3_workflow_users}
\renewcommand{\arraystretch}{1.35}
\rowcolors{2}{rowgray}{white}
\begin{tabularx}{\textwidth}{|L{4cm}|Y|}
\hline
\rowcolor{mainblue!15}
\textbf{Étape} & \textbf{Description} \\
\hline
Recherche et filtrage & L'administrateur retrouve un compte par nom, email, rôle ou statut. \\
\hline
Création ou modification & Le compte est créé ou mis à jour selon le besoin de gouvernance. \\
\hline
Changement de statut & Le statut actif, inactif ou suspendu est appliqué immédiatement. \\
\hline
Réinitialisation & Un accès temporaire est généré lorsque le mot de passe doit être réinitialisé. \\
\hline
Audit & Chaque action sensible alimente le journal d'audit. \\
\hline
\end{tabularx}
\end{table}

\paragraph{Interfaces réalisées}

Les écrans les plus parlants sont la \textbf{liste des comptes}, la \textbf{fiche de compte}
et le \textbf{dialogue de réinitialisation}. Ils permettent de comprendre vite le niveau
de contrôle offert à l'administrateur.

\bigscreen{0.92\textwidth}{images/admin_users.png}%
{Gestion avancée des utilisateurs}{fig:s3_ui_users}

\bigscreen{0.92\textwidth}{images/admin_user_detail.png}%
{Fiche de détail d'un utilisateur}{fig:s3_ui_user_detail}

\bigscreen{0.92\textwidth}{images/admin_user_reset.png}%
{Réinitialisation du mot de passe}{fig:s3_ui_user_reset}

\subsubsection{Module 4 — Messagerie temps réel}

La messagerie apporte de la fluidité dans le quotidien de la plateforme. Elle repose sur
des conversations persistées, des participants autorisés et un streaming SSE pour éviter
les rechargements inutiles.

\paragraph{Workflow du module}

\begin{table}[H]
\centering
\caption{Workflow du module Messagerie temps réel}
\label{tab:s3_workflow_chat}
\renewcommand{\arraystretch}{1.35}
\rowcolors{2}{rowgray}{white}
\begin{tabularx}{\textwidth}{|L{4cm}|Y|}
\hline
\rowcolor{mainblue!15}
\textbf{Étape} & \textbf{Description} \\
\hline
Création de conversation & L'utilisateur choisit les participants autorisés. \\
\hline
Ouverture du flux & Le frontend ouvre la connexion SSE pour la conversation courante. \\
\hline
Envoi d'un message & Un message est persisté dans l'historique de la conversation. \\
\hline
Diffusion temps réel & Le backend pousse immédiatement le message aux abonnés connectés. \\
\hline
Historique paginé & L'utilisateur retrouve les échanges sans saturer l'interface. \\
\hline
\end{tabularx}
\end{table}

\paragraph{Interfaces réalisées}

La présentation retenue repose sur trois vues complémentaires : la \textbf{liste des
conversations}, la \textbf{fenêtre de discussion} et le \textbf{panneau des participants}.
Ce choix permet de montrer la messagerie sans alourdir la lecture du chapitre.

\bigscreen{0.92\textwidth}{images/chat_messages.png}%
{Conversation et messages en temps réel}{fig:s3_ui_chat}

\bigscreen{0.92\textwidth}{images/chat_participants.png}%
{Participants d'une conversation}{fig:s3_ui_chat_participants}

\subsubsection{Module 5 — Gestion des membres d'équipe entreprise}

Ce module permet à une entreprise de piloter son équipe interne. Il reste volontairement
sobre : liste des membres, rôles, ajout, édition et suppression. C'est une fonctionnalité
de gouvernance, donc l'interface doit être efficace avant d'être décorative.

\paragraph{Workflow du module}

\begin{table}[H]
\centering
\caption{Workflow du module Gestion des membres d'équipe}
\label{tab:s3_workflow_team}
\renewcommand{\arraystretch}{1.35}
\rowcolors{2}{rowgray}{white}
\begin{tabularx}{\textwidth}{|L{4cm}|Y|}
\hline
\rowcolor{mainblue!15}
\textbf{Étape} & \textbf{Description} \\
\hline
Visualisation de l'équipe & L'entreprise consulte la liste de ses membres. \\
\hline
Ajout & Un nouveau membre est ajouté à l'espace entreprise. \\
\hline
Modification & Le rôle ou les coordonnées internes peuvent être ajustés. \\
\hline
Suppression & Un accès peut être retiré lorsqu'un membre quitte l'équipe. \\
\hline
Contrôle du périmètre & Les actions restent confinées à l'entreprise connectée. \\
\hline
\end{tabularx}
\end{table}

\paragraph{Interfaces réalisées}

Les interfaces sélectionnées sont la \textbf{liste d'équipe} et la \textbf{boîte d'édition
d'un membre}. Cela suffit pour montrer la logique sans surcharger le chapitre.

\bigscreen{0.92\textwidth}{images/team_members.png}%
{Gestion des membres d'équipe entreprise}{fig:s3_ui_team}

\bigscreen{0.92\textwidth}{images/team_member_form.png}%
{Formulaire d'ajout ou de modification d'un membre}{fig:s3_ui_team_form}

\subsubsection{Module 6 — Audit et traçabilité}

Le module d'audit est la pièce qui donne de la crédibilité à la plateforme. Il ne sert pas
à faire joli : il sert à prouver qui a fait quoi, quand, et sur quel objet.

\paragraph{Workflow du module}

\begin{table}[H]
\centering
\caption{Workflow du module Audit et traçabilité}
\label{tab:s3_workflow_audit}
\renewcommand{\arraystretch}{1.35}
\rowcolors{2}{rowgray}{white}
\begin{tabularx}{\textwidth}{|L{4cm}|Y|}
\hline
\rowcolor{mainblue!15}
\textbf{Étape} & \textbf{Description} \\
\hline
Déclenchement & Une action sensible est réalisée par un administrateur. \\
\hline
Lecture de l'état initial & Le backend récupère les anciennes valeurs avant modification. \\
\hline
Application de l'action & La modification métier est appliquée. \\
\hline
Enregistrement audit & Une trace détaillée est insérée dans le journal d'audit. \\
\hline
Consultation & L'historique est filtrable par type d'action et période. \\
\hline
\end{tabularx}
\end{table}

\paragraph{Interfaces réalisées}

Les vues retenues sont la \textbf{table d'audit filtrable}, la \textbf{fiche détaillée d'une
action} et le \textbf{résumé d'activité utilisateur}. Elles mettent en avant la valeur
du module de traçabilité.

\bigscreen{0.92\textwidth}{images/audit_history.png}%
{Historique d'audit et traçabilité}{fig:s3_ui_audit}

\bigscreen{0.92\textwidth}{images/audit_detail.png}%
{Détail d'une action d'audit}{fig:s3_ui_audit_detail}

\subsection{Tests}

\subsubsection{Tests fonctionnels}

Les tests fonctionnels ont été réalisés à l'aide de Postman pour le backend et par
validation manuelle des interfaces frontend. Ils couvrent les principaux parcours du
sprint et confirment le bon enchaînement entre les routes, les statuts métier et les
interfaces.

\begin{table}[H]
\centering
\caption{Scénarios de tests fonctionnels — Sprint~3}
\label{tab:s3_tests_fonctionnels}
\renewcommand{\arraystretch}{1.35}
\rowcolors{2}{rowgray}{white}
\begin{tabularx}{\textwidth}{|L{4cm}|Y|L{3.2cm}|}
\hline
\rowcolor{mainblue!15}
\textbf{Fonctionnalité} & \textbf{Cas testé} & \textbf{Résultat attendu} \\
\hline
Supervision & Accès au dashboard avec un compte inspecteur. & KPIs filtrés par territoire affichés. \\
\hline
Supervision & Accès au dashboard avec un compte ANETI. & Vision nationale complète affichée. \\
\hline
Reporting entreprise & Soumission d'un rapport de conformité. & Rapport créé avec statut \texttt{submitted}. \\
\hline
Reporting entreprise & Rejet sans recommandation. & Action bloquée, champ obligatoire. \\
\hline
Gestion utilisateurs & Changement de statut d'un compte. & Statut mis à jour et audit enregistré. \\
\hline
Gestion utilisateurs & Réinitialisation de mot de passe. & Accès temporaire généré et action tracée. \\
\hline
Messagerie & Envoi d'un message dans une conversation. & Message reçu en temps réel par l'autre participant. \\
\hline
Messagerie & Abonnement au flux SSE. & Nouveaux messages reçus sans rechargement. \\
\hline
Equipe entreprise & Ajout d'un membre d'équipe. & Membre enregistré pour l'entreprise connectée. \\
\hline
Audit & Consultation de l'historique d'un utilisateur. & Lignes d'audit filtrées et exploitables. \\
\hline
\end{tabularx}
\end{table}

\bigscreen{0.92\textwidth}{images/postman_sprint3.png}%
{Résultat du test Postman — endpoint de soumission d'un rapport de conformité}{fig:s3_postman}

\subsubsection{Tests d'intégration}

\begin{table}[H]
\centering
\caption{Scénarios de tests d'intégration — Sprint~3}
\label{tab:s3_tests_integration}
\renewcommand{\arraystretch}{1.35}
\rowcolors{2}{rowgray}{white}
\begin{tabularx}{\textwidth}{|L{5cm}|Y|L{3.2cm}|}
\hline
\rowcolor{mainblue!15}
\textbf{Interaction testée} & \textbf{Scénario} & \textbf{Résultat attendu} \\
\hline
Reporting entreprise $\rightarrow$ Supervision & Une entreprise soumet un rapport. & Le rapport apparaît dans la liste de supervision. \\
\hline
Supervision $\rightarrow$ Recommandations & Un inspecteur rejette un rapport avec recommandation. & Le statut devient \texttt{rejected} et la recommandation est persistée. \\
\hline
Gestion utilisateurs $\rightarrow$ Audit & Un admin modifie un utilisateur. & Une ligne d'audit est créée. \\
\hline
Messagerie $\rightarrow$ SSE & Un message est envoyé pendant qu'un participant écoute le flux. & Le message est diffusé au participant connecté. \\
\hline
RBAC $\rightarrow$ Routes protégées & Un utilisateur non autorisé tente d'accéder à une route admin. & Accès refusé HTTP~403. \\
\hline
\end{tabularx}
\end{table}

\section{Conclusion}

Le Sprint~3 ajoute à HandiTalents une véritable couche de pilotage et de preuve.
Les acteurs institutionnels disposent de tableaux de bord et de flux de validation, les
entreprises peuvent gérer leurs rapports de conformité, les utilisateurs échangent en
temps réel, l'administrateur contrôle les comptes et chaque action sensible est tracée.

Ce sprint est important parce qu'il donne au projet une maturité de plateforme :
il ne s'agit plus seulement de recruter, mais aussi de superviser, justifier, échanger
et auditer. Cette base prépare naturellement les évolutions ultérieures vers une
automatisation plus intelligente et une meilleure exploitation des données métiers.

\end{document}
